%Autor: miguev

% Documento de clase 'book', papel A4 a doble cara con letra de 12 puntos.
\documentclass[a4paper,10pt,twoside,final]{book}

% Paquetes
% ========

% Paquete pare que las fuentes se vean bien en los PDF. 
% Garcias a Jos� Miguel Espadero Guillermo por se�alarlo.
%\usepackage{pslatex}

% Paquete para poder escribir los caracteres especiales sin necesidad de
% comandos. La opci�n 'latin1' se corresponde con el standar ISO-8859-1
% (tambi�n es v�lido para ISO-8859-15, aunque no reconoce el �)
\usepackage[latin1]{inputenc}

% Paquete para castellanizar las variables del documento.
\usepackage[spanish,activeacute]{babel}

% Paquete de s�mbolos de LaTeX y de la AMS (American Mathematical Society)
\usepackage{latexsym}
\usepackage{amssymb}

% Estilo de cabeceras y pi�s de p�gina.
\usepackage{fancyhdr}

% Paquetes para inclusi�n de gr�ficos Encapsulados PostScript
\usepackage{graphicx}
%\usepackage{psfig}
\usepackage{subfigure}

% Paquete para el s�mbolo oficial del �uro
\usepackage{eurosym}

% Paquete para construir los �ndices con makeindex.
\usepackage{makeidx}

% Paquete para hiperenlaces HTML
%\usepackage[colorlinks=true,urlcolor=blue]{hyperref}

% Paquete para filas m�ltiples (para tablas)
\usepackage{multirow}

% Paquete para usar colores, en las p�ginas y en las tablas
\usepackage{color}
\usepackage{colortbl}

% Paquete para vieguer�as con verbatim
\usepackage{verbatim}

% Paquete para las marcas de corte, para imprenta
\usepackage{crop}

% Estilo para los apuntes del TILA, 
% tomando c�digo de los paquetes fncychap y ull.
% Este paquete debe ser el �ltimo en ser incluido, 
% ya que sobreescribe definiciones de otros paquetes
% (p.ej. babel)
\usepackage{cila}

\frenchspacing

% Habilida las instrucciones especiales de indexado
\makeindex

% COMIENZO DEL DOCUMENTO
\begin{document}

% Se da la orden de poner las marcas de corte
\crop[cam,info,horigin=-0.41in,vorigin=-5mm]

% El  comando  \sloppy relaja  los  p�rrafos  y  evita muchos  de  los
% desbordamientos  horizontales. Esto  se  logra  permitiendo que  los
% espacios entre palabras lleguen a ser m�s grandes.
\sloppy

% Parte frontal
% Portadas, nota de copyright y dedicatoria
\frontmatter
%Autor: miguev

\pagestyle{empty}

\vspace*{20mm}
\begin{center}
{\LARGE Gu�a de aprendizaje para estudiantes}\vspace*{5mm}
%\resizebox{\textwidth}{10mm}{{\LARGE\bf Curso de Introducci�n a Linux para Alumnos}}
%\resizebox{\textwidth}{10mm}{{\LARGE\bf Introducci�n a Linux y el Software Libre}}
\resizebox{\textwidth}{10mm}{{\LARGE\bf \titulo}}
\end{center}

\vfill
\cleardoublepage

\vspace*{20mm}
\begin{center}
{\LARGE Gu�a de aprendizaje para estudiantes}\vspace*{5mm}
%\resizebox{\textwidth}{10mm}{{\LARGE\bf Curso de Introducci�n a Linux para Alumnos}}
%\resizebox{\textwidth}{10mm}{{\LARGE\bf Introducci�n a Linux y el Software Libre}}
\resizebox{\textwidth}{10mm}{{\LARGE\bf \titulo}}
\end{center}

\begin{center}
\vspace*{60mm}
\includegraphics[width=120mm]{imagenes/gulic.eps}
\end{center}

\vfill

\begin{flushright}
%{\LARGE Grupo de Usuarios de Linux de Canarias}
%\vspace*{5mm}
{\large
Miguel �ngel Vilela Garc�a     \\ % 69 pags
Jes�s Miguel Torres Jorge      \\ % 65 pags
Carlos P�rez P�rez             \\ % 51 pags
Tom�s Bautista Delgado         \\ % 50 pags
Carlos Alberto Morales D�az    \\ % 44 pags
F�lix J. Marcelo Wirnitzer     \\ % 30 pags
Carlos de la Cruz Pinto        \\ % 18 pags
Ren� Mart�n Rodr�guez          \\ % 16 pags
Ed�n Kozo                      \\ % 13 pags
Carlos Mestre Gonz�lez         \\ % 11 pags
% Teresa Gonz�lez de la F�       \\ %  9 pags
% Luis Cabrera Sa�co             \\ %  8 pags
% Pedro Gracia Fajardo           \\ %  2 pags
% Carlos A. Paramio              \\ %  1 pags
}
\end{flushright}

% Autores ordenados por cantidad de p�ginas:
% Miguel �ngel Vilela Garc�a
% Jes�s Miguel Torres Jorge
% Carlos P�rez P�rez
% Tom�s Bautista Delgado
% Carlos Alberto Morales D�az
% F�lix J. Marcelo Wirnitzer
% Carlos de la Cruz Pinto
% Ren� Mart�n Rodr�guez
% Ed�n Kozo
% Carlos Mestre Gonz�lez
% Teresa Gonz�lez de la F�
% Luis Cabrera Sa�co
% Pedro Gracia Fajardo
% Carlos A. Paramio

\vfill
\clearpage

\vspace*{130mm}
\noindent Copyright  \copyright~ 2001 -- 2003 \GULiC
\vspace*{5mm}

\noindent  Se da  permiso  para copiar,  distribuir  o modificar  este
documento  en los  t�rminos que  establece la  GNU Free  Documentation
License, Versi�n 1.1 o cualquier  otra versi�n posterior publicada por
la Free Software  Foundation; las secciones que  no pueden modificarse
son  ``Presentaci�n''  y  ``Agradecimientos'',  y  no  existen  textos
cubierta ni contracubierta. Se incluye una  copia de la licencia en la
secci�n ``GNU Free Documentation  License''. Este libro es \copyright~
de sus autores.

\noindent  La forma  original de  este documento  es c�digo  fuente en
\LaTeX.  La compilaci�n  de este  c�digo  fuente en  \LaTeX~ tiene  el
efecto de  generar una representaci�n del  documento independiente del
dispositivo, que puede convertirse a otros formatos o imprimirse.

\noindent  La  GNU  Free  Documentation  License  est�  disponible  en
www.gnu.org o solici�ndola por escrito a the Free Software Foundation,
Inc., 59 Temple Place - Suite 330, Boston, MA 02111-1307, USA.

\noindent  Este  documento lo  han  maquetado  sus autores  utilizando
\LaTeX~ y {\sf The GIMP}, que son programas abiertos y libres.

\vfill
\clearpage

\vspace*{35mm}
\begin{flushright}
{\large A los novatos y novatas en el mundo de GNU/Linux, \\
y a sus familiares, amistades y parejas} \\
{\small (ellas y ellos saben por qu�)}
\end{flushright}

\vfill
\cleardoublepage

\pagestyle{fancy}


% Presentaci�n y Contribuciones
\pagestyle{plain}
\incluye{presentacion}
\incluye{contrib}
\newpage
\pagestyle{fancy}

% Modificaci�n de algunas variables para adaptarlas a nuestro uso particular
\renewcommand{\partname}{M�DULO}
\renewcommand{\chaptername}{TEMA}
\renewcommand{\listtablename}{�ndice de tablas}
\renewcommand{\tablename}{Tabla}

% �ndices general, de tablas, de figuras y de ejemplos
\tableofcontents\newpage\thispagestyle{empty}
\listoffigures\newpage\thispagestyle{empty}
\listoftables\newpage\thispagestyle{empty}
\listadeejemplos\newpage\thispagestyle{empty}

% Parte principal
\mainmatter

% Inclusi�n de los ficheros en el orden apropiado.
%\incluye{administracion}

% M�dulos
% I.   Entorno GNU/Linux (3 d�as, 15 horas en clases de 5 horas)
% II.  Party (2 d�as, 20 horas en jornadas de 10 horas)
% III. Documentaci�n (5 d�as, 20 horas, en clases de 4 horas)
% IV.  Matem�ticas (3 d�as, 12 horas, en clases de 4 horas)
% V.   Programaci�n (5 d�as, 20 horas, en clases de 4 horas)
% VI.  Programaci�n avanzada (?? horas)

% M�dulo I:   Entorno GNU/Linux
% 1. Historia de GNU/Linux
% 2. El entorno X-Window (Gnome, KDE)
% 3. El entorno del int�rprete de comandos (con mtools)
% 4. Usando internet (mozilla, konqueror, kmail, gftp, mutt, SSH)
% 5. Aplicaciones (gnotepad, abiword, openoffice)
% 6. Documentaci�n y sistemas de ayuda
\part{Entorno GNU/Linux}
\incluye{introduccion}
\incluye{comandos}
\incluye{xwindow}
\incluye{internet}
%\incluye{aplicaciones}
%\incluye{documentacion}

% M�dulo II:  Party
% 1. Sistema base (particiones)
% 2. Hardware y kernel (make [menu|x]config && make-kpkg)
% 3. Software para los m�dulos.
% 4. Administraci�n b�sica (adduser, floppy, cdrom, APT)
% 5. Paquetes en fuentes (bajar, compilar, instalar Xine+dvdnav)
% 6. Software adicional (a gusto del consumidor)
%\part{Instalaci�n de GNU/Linux}
%\incluye{instalacion.tex}

% M�dulo III: Edici�n y maquetaci�n de documentos
% 1. Edici�n de gr�ficos (DIA, QCad, The GIMP)
% 2. OpenOffice
% 3. HTML (lenguaje, bluefish, quanta)
% 4. LyX
% 5. LaTeX
\part{Edici�n de gr�ficos y documentos}
\incluye{graficos}
\incluye{openoffice}
\incluye{html}
%\incluye{lyx}
%\incluye{latex}

% M�dulo IV:  Matem�ticas
% 1. GNUplot
% 2. Octave (SciLab)
% 3. R
% 4. Yacas
% 5. Maxima
\part{Herramientas matem�ticas}
\incluye{octave}
\incluye{gnuplot}
\incluye{r}
\incluye{yacas}

% M�dulo V:   Programaci�n
% 1. Editores de texto (gvim, xemacs)
% 2. Bash
% 3. GNU Fortran 77
% 4. FreePascal
% 5. GNU C/C++
% 6. Java
% 7. GNU Make
% 8. Depuradores (gdb y ddd)
\part{Programaci�n}
\incluye{editores}
\incluye{freepascal}
\incluye{gnufortran}
\incluye{gnuc}
\incluye{gnumake}
\incluye{depuradores}
\incluye{java}
\incluye{bash}

%
% NO SE ENTRETENGAN CON ESTO TODAV�A
%
% M�dulo VI:  Programaci�n avanzada
% 1. PHP + [My|Postgre]SQL
% 2. Perl
% 3. Python (NumPy) 
% 4. GUIs toolkits: Tk, QT, Wx, Gtk, Glade
% 5. XML
%\part{Programaci�n avanzada}
%\incluye{php}
%\incluye{perl}
%\incluye{python}
%\incluye{toolkits}
%\incluye{xml}
%\indluye{ides}


% Ap�ndices y referencias
\appendix
\incluye{recursos}
%\incluye{fdles}




% Parte posterior
\backmatter

% �ndice terminol�gico
\printindex\newpage\thispagestyle{empty}

% Bibliograf�a, generado por BIBTeX
\bibliographystyle{plain}
\bibliography{biblio}

\end{document}
% FIN DEL DOCUMENTO
